\documentclass[11pt]{article}

% ---------- Packages ----------
\usepackage[a4paper, margin=2.4cm]{geometry}
\usepackage[utf8]{inputenc}
\usepackage[T1]{fontenc}
\usepackage{lmodern}
\usepackage[spanish]{babel}
\usepackage{xcolor}
\usepackage{titlesec}
\usepackage{enumitem}
\usepackage{fancyhdr}
\usepackage{hyperref}
\usepackage{tikz}
\usetikzlibrary{positioning, arrows.meta, fit, backgrounds, shapes.misc}
\usepackage{tcolorbox}
\tcbuselibrary{skins, breakable}

% ---------- Colors ----------
\definecolor{inkNavy}{HTML}{1B2A4A}
\definecolor{accentTeal}{HTML}{0E7C7B}
\definecolor{accentAmber}{HTML}{C97A2C}
\definecolor{softGray}{HTML}{6B7280}
\definecolor{bgPlanner}{HTML}{EAF3F3}
\definecolor{bgBuild}{HTML}{FBF1E6}
\definecolor{bgAudit}{HTML}{EEF0F8}
\definecolor{lineGray}{HTML}{D7DBE0}

% ---------- Heading typography ----------
\titleformat{\section}
  {\Large\bfseries\color{inkNavy}}
  {\thesection}{0.5em}{}
\titlespacing*{\section}{0pt}{1.4em}{0.8em}

\titleformat{\subsection}
  {\large\bfseries\color{accentTeal}}
  {\thesubsection}{0.5em}{}
\titlespacing*{\subsection}{0pt}{1.1em}{0.5em}

% ---------- Header/footer ----------
\pagestyle{fancy}
\fancyhf{}
\renewcommand{\headrulewidth}{0pt}
\fancyfoot[L]{\small\color{softGray}\href{https://github.com/dramirezbe/SelfHostedContextCrafter}{SelfHostedContextCrafter} --- Documento conceptual}
\fancyfoot[R]{\small\color{softGray}\thepage}

\hypersetup{
  colorlinks=true,
  linkcolor=accentTeal,
  urlcolor=accentTeal,
  pdftitle={SelfHostedContextCrafter},
  pdfauthor={SelfHostedContextCrafter},
  pdfsubject={Orquestador de agentes autohospedado para planificación acotada y auditoría de código}
}

\setlist[itemize]{leftmargin=1.4em, itemsep=0.3em, topsep=0.3em}

% Reusable concept box
\newtcolorbox{ideabox}[1]{
  enhanced, breakable,
  colback=white, colframe=lineGray,
  boxrule=0.5pt, left=10pt, right=10pt, top=8pt, bottom=8pt,
  borderline west={2.5pt}{0pt}{accentTeal},
  before upper={\textbf{\color{inkNavy}#1}\\[4pt]}
}

\begin{document}

% ================= COVER PAGE =================
\begin{titlepage}
\centering
\vspace*{2.2cm}

{\color{softGray}\large ORQUESTADOR DE AGENTES AUTOHOSPEDADO}\\[0.6cm]

{\Huge\bfseries\color{inkNavy} \href{https://github.com/dramirezbe/SelfHostedContextCrafter}{SelfHostedContextCrafter}}\\[0.5cm]

{\Large\color{accentTeal} Planificación acotada y auditoría de código\\ sobre infraestructura local}

\vspace{1.8cm}

\begin{tikzpicture}
\draw[accentAmber, line width=1.2pt] (-3,0) -- (3,0);
\end{tikzpicture}

\vspace{1.8cm}

{\large\color{softGray} Documento conceptual --- Propuesta de idea}

\vspace{0.4cm}

{\color{softGray} \today}

\vfill

{\small\color{softGray} Un sistema que separa \emph{pensar} (planificar con contexto verificable),\\ \emph{construir} (hacer con un humano al control) y \emph{verificar} (auditar contra la fuente),\\ todo ejecutándose en tu propio hardware.}

\end{titlepage}

% ================= EXECUTIVE SUMMARY =================
\section*{Resumen ejecutivo}
\addcontentsline{toc}{section}{Resumen ejecutivo}

\href{https://github.com/dramirezbe/SelfHostedContextCrafter}{SelfHostedContextCrafter} es un orquestador de agentes LLM \textbf{autohospedado} diseñado para equipos que desarrollan software bajo restricciones del mundo real --- hardware específico, requisitos de negocio y normativas --- que no pueden (o no quieren) enviar esa información a un proveedor de IA en la nube.

La idea central es simple pero poco común entre las herramientas actuales: \textbf{el agente nunca escribe código de producción}. Su única tarea es leer y digerir toda la documentación técnica relevante (hojas de datos de hardware, requisitos, normativas, contexto de arquitectura) y producir un \emph{plan avanzado} acotado y trazable. Ese plan es lo que un humano --- apoyado por su propio agente de código --- usa para construir. Una vez que el código existe, un segundo modo del mismo sistema lo \textbf{audita} contra el plan y la documentación original, devolviendo una puntuación por módulo, cambios sugeridos y advertencias específicas.

Todo el ciclo se ejecuta en hardware local, con modelos LLM autohospedados, lo que hace viable este enfoque en contextos donde la confidencialidad del hardware, código o normativas es un requisito, no una preferencia.

% ================= THE PROBLEM =================
\section{El problema}

Los asistentes de código actuales son potentes generando texto, pero comparten tres debilidades recurrentes en proyectos con restricciones del mundo real:

\begin{itemize}
  \item \textbf{Pierden el contexto original.} Generan código con apariencia razonable en el momento, pero sin trazabilidad clara hacia el requisito, normativa u hoja de datos de hardware que lo originó.
  \item \textbf{Alucinan cuando el contexto está disperso.} Cuanto más fragmentada está la documentación (PDFs del fabricante, normativas, especificaciones internas), más fácil es que el modelo rellene vacíos con suposiciones.
  \item \textbf{Requieren enviar información sensible a servicios externos.} El hardware propietario, las normativas internas o el código fuente no siempre pueden salir de la organización.
\end{itemize}

A esto se suma un problema de \textbf{control}: cuando un agente genera y audita su propio código en el mismo paso, no hay verificación independiente real --- es el mismo razonamiento revisándose a sí mismo.

% ================= THE CORE IDEA =================
\section{La idea central}

\href{https://github.com/dramirezbe/SelfHostedContextCrafter}{SelfHostedContextCrafter} divide el ciclo de desarrollo en tres roles distintos, cada uno con una responsabilidad acotada:

\begin{ideabox}{1. Planificar (Modo Planner)}
Ingiere toda la documentación técnica disponible y la convierte en un plan de proyecto avanzado. No escribe una sola línea de código de producción: su entregable es la estructura del proyecto y un plan acotado a lo que la documentación realmente dice.
\end{ideabox}

\vspace{0.4em}

\begin{ideabox}{2. Construir (Humano + agente de código)}
Un desarrollador, apoyado por su agente de código de preferencia, ejecuta el plan. Aquí es donde se escribe realmente el software --- con criterio humano, y con el plan acotado como guía, no como camisa de fuerza.
\end{ideabox}

\vspace{0.4em}

\begin{ideabox}{3. Auditar (Modo Auditor)}
El mismo sistema, conservando la memoria de todo lo ingerido durante la fase de planificación, revisa el código resultante contra el plan original y la documentación fuente. Devuelve una puntuación por módulo, cambios sugeridos y advertencias específicas.
\end{ideabox}

\vspace{0.6em}

La clave no es ninguno de los tres pasos por sí solo, sino que los tres comparten la \textbf{misma base de conocimiento retenida} (RAG). El auditor no audita en el vacío: audita contra exactamente el mismo contexto normativo y técnico que generó el plan, lo que hace que sus advertencias sean trazables hasta la fuente original.

% ================= ARCHITECTURE =================
\section{Arquitectura de flujo conceptual}

\begin{figure}[h!]
\centering
\resizebox{\textwidth}{!}{%
\begin{tikzpicture}[
  font=\small,
  box/.style={draw, rounded corners=2pt, align=center, minimum height=0.9cm, inner sep=6pt, line width=0.6pt},
  input/.style={box, draw=accentTeal, fill=white, text width=2.6cm},
  proc/.style={box, draw=inkNavy, fill=inkNavy!85, text=white, text width=3.2cm, font=\small\bfseries},
  outbox/.style={box, draw=accentTeal!70, fill=white, text width=2.6cm},
  human/.style={box, draw=accentAmber, fill=accentAmber!12, text width=4.6cm, minimum height=1.1cm},
  note/.style={font=\scriptsize\itshape, text=softGray, align=center},
  arr/.style={-{Stealth[length=2.2mm]}, line width=0.7pt, color=softGray},
  stagebg/.style={rounded corners=4pt, line width=0pt},
  stagelabel/.style={font=\bfseries\small, text=inkNavy}
]

% ---------- COLUMN 1: PLANNER MODE ----------
\node[input] (hw) at (0,8.0) {Contexto HW};
\node[input] (req) at (0,6.3) {Contexto de requisitos};
\node[input] (norm) at (0,4.6) {Contexto normativo};
\node[input, text width=2.9cm] (esp) at (0,2.3) {Contexto específico\\(arq./protocolos/\\advertencias)};

\node[proc, text width=2.6cm] (conv1) at (3.6,5.2) {Conversor\\PDF \(\to\) MD};
\node[proc, text width=2.9cm, minimum height=1.6cm] (fw1) at (7.4,5.2) {FRAMEWORK\\\footnotesize (Modo Planner)};
\node[note, text width=2.9cm] at (7.4,4.0) {retiene RAGs,\\PDFs y MDs};

\node[outbox] (struct) at (11.0,6.6) {Estructura del proyecto};
\node[outbox] (plan) at (11.0,3.8) {Plan avanzado};

\draw[arr] (hw) -- (conv1);
\draw[arr] (req) -- (conv1);
\draw[arr] (norm) -- (conv1);
\draw[arr] (esp) -- (conv1);
\draw[arr] (conv1) -- (fw1);
\draw[arr] (fw1) -- (struct);
\draw[arr] (fw1) -- (plan);

\begin{scope}[on background layer]
\node[stagebg, fill=bgPlanner, fit=(hw)(req)(norm)(esp)(conv1)(fw1)(struct)(plan), inner sep=10pt] (stage1) {};
\end{scope}
\node[stagelabel] at (stage1.north) [above, yshift=2pt] {MODO 1 --- PLANNER};

% ---------- STAGE 2: HUMAN + AGENT ----------
\node[human, text width=5.8cm] (build) at (7.4,-0.6) {Humano + agente de código\\\footnotesize construye el software};
\node[outbox, fill=white] (code) at (7.4,-2.3) {Código fuente};

% Advanced plan -> Build (right-angle route, avoiding text overlap)
\draw[arr] (plan.south) -- (11.0,2.0) -- (7.4,2.0) -- (build.north);
\draw[arr] (build) -- (code);

\begin{scope}[on background layer]
\node[stagebg, fill=bgBuild, fit=(build), inner sep=10pt] (stage2) {};
\end{scope}
\node[stagelabel] at (stage2.north) [above, yshift=2pt] {CONSTRUIR};

% ---------- COLUMN 3: AUDITER MODE ----------
\node[proc, text width=2.6cm] (conv2) at (3.6,-4.5) {Conversor\\(análisis de código)};
\node[proc, text width=2.9cm, minimum height=1.6cm] (fw2) at (7.4,-5.2) {FRAMEWORK\\\footnotesize (Modo Auditor)};
\node[note, text width=3.1cm] at (7.4,-6.4) {retiene RAGs\\y agentes\\(continuidad del Modo 1)};

\node[outbox] (score) at (11.0,-3.9) {Puntuación por módulo};
\node[outbox] (changes) at (11.0,-5.2) {Cambios sugeridos};
\node[outbox] (warn) at (11.0,-6.5) {Advertencias específicas};

\draw[arr] (code) -- (conv2);
\draw[arr] (conv2) -- (fw2);
\draw[arr] (fw2) -- (score);
\draw[arr] (fw2) -- (changes);
\draw[arr] (fw2) -- (warn);

\begin{scope}[on background layer]
\node[stagebg, fill=bgAudit, fit=(conv2)(fw2)(score)(changes)(warn), inner sep=10pt] (stage3) {};
\end{scope}
\node[stagelabel] at (stage3.south) [below, yshift=-2pt] {MODO 2 --- AUDITOR};

% ---------- FEEDBACK LOOP (outer route, no crossings) ----------
\draw[arr, color=accentAmber, line width=1pt] (changes.east) -- (13.2,-5.2) -- (13.2,-0.6) -- (build.east);
\node[note, text width=2.6cm, text=accentAmber] at (13.2,-2.9) [right] {itera la\\construcción};

\end{tikzpicture}%
}
\caption{\small Flujo conceptual: el contexto fluye desde el Planner hacia la etapa de Construcción y de allí al Auditor, sin que el agente toque código de producción directamente. Los cambios sugeridos y las advertencias se realimentan al ciclo de construcción.}
\end{figure}

\subsection{Leyendo el diagrama}

\begin{itemize}
  \item \textbf{El Modo 1} normaliza cuatro fuentes de contexto distintas (hardware, requisitos, normativas, contexto específico) a través de un único conversor PDF\,\(\to\)\,MD, para que el Framework trabaje siempre sobre texto homogéneo. La salida no es código: es una estructura de proyecto y un plan avanzado, y el Framework retiene el material fuente como una base de conocimiento persistente.
  \item \textbf{La etapa de construcción} es deliberadamente humana. El plan avanzado y el contexto retenido alimentan al desarrollador y su agente de código de preferencia, pero las decisiones de implementación permanecen en manos humanas.
  \item \textbf{El Modo 2} recibe el código resultante (a través de un conversor dedicado orientado al análisis de código) y el mismo plan avanzado. El Framework en este modo no parte de cero: retiene los RAGs y agentes del Modo 1, por lo que puede auditar el código contra la documentación original, no solo contra el plan.
\end{itemize}

% ================= DESIGN PRINCIPLES =================
\section{Principios de diseño}

\begin{ideabox}{Autohospedado por defecto}
Modelos LLM ejecutándose en hardware local. Ninguna hoja de datos de hardware, normativa interna o línea de código necesita salir de tu propia infraestructura para que el sistema funcione.
\end{ideabox}

\vspace{0.4em}

\begin{ideabox}{El agente planifica, el humano construye}
Separar la generación del plan de la generación de código reduce la superficie de alucinación donde más duele: las decisiones arquitectónicas. El código lo sigue escribiendo quien tiene el criterio y la responsabilidad final.
\end{ideabox}

\vspace{0.4em}

\begin{ideabox}{Planificación acotada, no plantillas genéricas}
El plan avanzado no proviene del conocimiento general del modelo: proviene de lo que la documentación ingerida realmente dice. Eso lo hace más acotado, pero también más fiable y más fácil de justificar en una auditoría real.
\end{ideabox}

\vspace{0.4em}

\begin{ideabox}{Auditoría con memoria, no auditoría ciega}
Al retener RAGs y agentes entre modos, el auditor evalúa el código contra la misma fuente de verdad que generó el plan --- no contra una interpretación nueva y potencialmente diferente del mismo problema.
\end{ideabox}

% ================= USE CASES =================
\section{Dónde encaja esta idea}

\begin{itemize}
  \item Equipos de \textbf{firmware o sistemas embebidos} que trabajan con hojas de datos de fabricantes y normativas de seguridad (industrial, automotriz, médica) y que no pueden subir esos documentos a un LLM en la nube.
  \item Organizaciones con \textbf{requisitos de soberanía de datos} que ya tienen, o están dispuestas a invertir en, cómputo local para inferencia de modelos.
  \item Proyectos donde el \textbf{costo de una mala interpretación normativa} es alto, y donde tener un plan trazable a la documentación fuente tiene valor por sí mismo, más allá de la velocidad de desarrollo.
\end{itemize}

% ================= CLOSING =================
\section{En una frase}

\begin{center}
\begin{tcolorbox}[
  enhanced, colback=inkNavy, colframe=inkNavy,
  coltext=white, width=0.9\textwidth, halign=center,
  arc=3pt, boxrule=0pt
]
\large\itshape ``Un orquestador que separa pensar, construir y verificar --- uno que nunca asume lo que la documentación no dice, y nunca abandona el hardware sobre el que se ejecuta.''
\end{tcolorbox}
\end{center}

\end{document}
